%% ----------------------------------------------------------------------------
%%  §3.10  Appendice: successioni estratte, Bolzano–Weierstrass, Cauchy
%% ----------------------------------------------------------------------------
\section[Appendice: successioni estratte, teorema di Bolzano--Weierstrass, successioni di Cauchy]%
  {Appendice: successioni estratte, teorema di\\ Bolzano--Weierstrass, successioni di Cauchy}
\label{sec:bw-cauchy}

\subsection{Successioni estratte}

\begin{definizione}[Successione estratta]\index{successione!estratta}
Sia $a_{n}$ una successione di numeri reali e sia $n_{k}$ una successione strettamente
crescente di numeri naturali. La successione $a_{n_{k}}$ definita da
\begin{equation}\label{eq:estratta}
  k\in\N\ \longmapsto\ a_{n_{k}}
\end{equation}
prende il nome di \emph{successione estratta} (o \emph{sottosuccessione}) da $a_{n}$, di
indici $n_{k}$.
\end{definizione}

In altre parole, un'estratta si ottiene «saltando» alcuni termini della successione (anche
infiniti), ma conservando l'ordine. Ad esempio, se $n_{k}=2k$, la successione estratta da
$a_{n}$ di indici $2k$, cioè di indici pari, è
\begin{equation}\label{eq:estratta-pari}
  a_{2},\ a_{4},\ \dots,\ a_{2k},\ \dots
\end{equation}
Se invece $n_{k}=2k-1$, si ottiene l'estratta da $a_{n}$ di indici dispari:
\begin{equation}\label{eq:estratta-dispari}
  a_{1},\ a_{3},\ \dots,\ a_{2k-1},\ \dots
\end{equation}

\begin{lemma}\label{lem:nk-maggiore-k}
Per ogni successione $n_{k}$ strettamente crescente di numeri naturali si ha
\begin{equation}\label{eq:nk-k}
  n_{k}\ge k \qquad\perogni k\in\N .
\end{equation}
\end{lemma}

\begin{schema}
\Ipotesi{$n_{k}\in\N$, \ $n_{k}<n_{k+1}$ per ogni $k$.}
\Tesi{$n_{k}\ge k$ per ogni $k\in\N$.}
\end{schema}

\begin{strategia}
Induzione su $k$: il primo indice è almeno $1$, e a ogni passo l'indice cresce almeno di
un'unità, perché tra due naturali consecutivi non ci sono altri naturali.
\end{strategia}

\begin{dimostrazione}
Per $k=1$ si ha ovviamente $n_{1}\ge1$. Inoltre, supponendo valida la \eqref{eq:nk-k},
proviamo che risulta $n_{k+1}\ge k+1$, da cui, per il principio di induzione, la
\eqref{eq:nk-k} risulterà vera per ogni $k$. Per ipotesi è $n_{k+1}>n_{k}\ge k$, ovvero
$n_{k+1}>k$, e perciò $n_{k+1}\ge k+1$ (\cref{lem:struttura-N}(iii)).
\end{dimostrazione}

Dalla \eqref{eq:nk-k} si ricava facilmente la seguente

\begin{proposizione}\label{prop:estratta-limite}
Se $a_{n}$ ha limite $a$ (finito o infinito), allora ogni estratta $a_{n_{k}}$ ha lo
stesso limite $a$.
\end{proposizione}

\begin{schema}
\Ipotesi{$a_{n}\to a$, con $a\in\R$ oppure $a=\pm\infty$; \ $n_{k}$ strettamente
crescente.}
\Tesi{$a_{n_{k}}\to a$ per $k\to+\infty$.}
\end{schema}

\begin{strategia}
L'estratta prende i suoi termini sempre più avanti nella successione: per il lemma
precedente il $k$-esimo termine dell'estratta ha indice $n_{k}\ge k$. Quindi la soglia
$\nu$ che funziona per la successione funziona anche per l'estratta.
\end{strategia}

\begin{dimostrazione}
Sia $a$ finito. Fissato $\epsilon>0$ esiste $k_{0}$ tale che $\abs{a_{n}-a}<\epsilon$ per
ogni $n>k_{0}$. Se $k>k_{0}$, essendo $n_{k}\ge k$ per il lemma precedente, si ha anche
$n_{k}>k_{0}$, e perciò $\abs{a_{n_{k}}-a}<\epsilon$. Se $a=\pm\infty$ si ragiona allo
stesso modo, con la disuguaglianza $a_{n}>M$ (o $a_{n}<-M$) al posto di
$\abs{a_{n}-a}<\epsilon$.
\end{dimostrazione}

\begin{corollario}\label{cor:estratte-limiti-diversi}
Se una successione ha due estratte con limiti diversi, allora non ha limite.
\end{corollario}

\begin{schema}
\Ipotesi{$a_{n_{k}}\to\ell_{1}$, \ $a_{m_{k}}\to\ell_{2}$ (finiti o infiniti), con
$\ell_{1}\neq\ell_{2}$.}
\Tesi{$a_{n}$ non ha limite, né finito né infinito.}
\end{schema}

\begin{strategia}
È la \cref{prop:estratta-limite} letta «al contrario»: se la successione avesse un
limite, tutte le estratte avrebbero quello stesso limite, e per l'unicità del limite non
potrebbero averne due diversi.
\end{strategia}

\begin{dimostrazione}
Se fosse $a_{n}\to a$ (finito o infinito), per la \cref{prop:estratta-limite} si avrebbe
$a_{n_{k}}\to a$ e $a_{m_{k}}\to a$. Per l'unicità del limite (\cref{teo:unicita-limite}
e, per i limiti infiniti, l'osservazione che segue la definizione di successione regolare
nel \cref{sec:def-limite}) sarebbe $\ell_{1}=a=\ell_{2}$, contro l'ipotesi.
\end{dimostrazione}

Ad esempio, $(-1)^{n}$ non ha limite perché l'estratta di indici pari tende a $1$ e quella
di indici dispari tende a $-1$: è una dimostrazione della
\cref{prop:alterna-non-converge} più rapida di quella data a suo tempo. Il
\cref{teo:lim-geometrica} per $a<-1$ si ottiene allo stesso modo.

\subsection{Il teorema di Bolzano--Weierstrass}

Nel \cref{sec:operazioni-limiti} abbiamo dimostrato che ogni successione convergente è
limitata, cioè esiste $M>0$ tale che $\abs{a_{n}}\le M$ per ogni $n\in\N$. Il viceversa non
sussiste, perché, ad esempio, la successione $a_{n}=(-1)^{n}$ è limitata ma non
convergente. Tuttavia, in essa si possono individuare estratte convergenti (quella dei
termini di posto pari, ad esempio). Ciò non è un caso: sussiste un notevole teorema,
che nello schema dei grandi teoremi di esistenza della \cref{fig:schema-teoremi} segue
immediatamente quello sulle successioni monotòne.

Per enunciarlo e dimostrarlo conviene precisare il linguaggio. Diciamo che un intervallo
$I$ \emph{contiene termini della successione per infiniti indici} se l'insieme
$\{n\in\N : a_{n}\in I\}$ non è limitato superiormente, cioè se per ogni $m\in\N$ esiste
$n>m$ con $a_{n}\in I$. Se $I=I'\cup I''$ e né $I'$ né $I''$ contengono termini per
infiniti indici, allora gli insiemi $\{n : a_{n}\in I'\}$ e $\{n : a_{n}\in I''\}$ sono
limitati superiormente, e quindi lo è anche la loro unione $\{n : a_{n}\in I\}$:
\emph{se $I$ contiene termini per infiniti indici, almeno una delle due parti $I'$, $I''$
ne contiene per infiniti indici}.

\begin{teorema}[di Bolzano--Weierstrass]\label{teo:bolzano-weierstrass}\index{teorema!di Bolzano--Weierstrass}
Sia $a_{n}$ una successione limitata. Allora esiste almeno una sua estratta convergente.
\end{teorema}

\begin{schema}
\Ipotesi{$A\le a_{n}\le B$ per ogni $n\in\N$.}
\Tesi{$\esiste n_{1}<n_{2}<\cdots<n_{k}<\cdots$ e $\ell\in[A,B]$ tali che $a_{n_{k}}\to\ell$.}
\end{schema}

\begin{strategia}
È il metodo di bisezione, già incontrato per approssimare $\sqrt2$, usato questa volta per
\emph{catturare} un valore attorno al quale si addensano infiniti termini della
successione. La successione ha infiniti termini, e sono
tutti nell'intervallo $[A,B]$. Dividiamo l'intervallo a metà (\cref{fig:bolzano-weierstrass}): almeno una delle due metà
contiene termini della successione per infiniti indici (se ciascuna ne contenesse solo un
numero finito, i termini sarebbero in tutto un numero finito). Scegliamo quella metà e
ripetiamo: otteniamo una successione di intervalli incapsulati $[A_{k},B_{k}]$, ciascuno
lungo la metà del precedente e ciascuno contenente termini della successione per infiniti
indici. Gli estremi sinistri $A_{k}$ formano una successione crescente e limitata, che
converge (\cref{teo:monotone}) a un numero $\ell$: qui entra in gioco la completezza.
Infine scegliamo un termine $a_{n_{k}}$ in ciascun intervallo, con indici crescenti
(possibile proprio perché ogni intervallo ne contiene per infiniti indici): l'estratta
così ottenuta è stretta tra $A_{k}$ e $B_{k}$, che tendono entrambi a $\ell$, e i
carabinieri concludono.
\end{strategia}

\begin{dimostrazione}
Per ipotesi la successione $a_{n}$ è limitata; pertanto esistono due costanti $A,B\in\R$,
con $A<B$, tali che
\begin{equation}\label{eq:bw-limitata}
  A\le a_{n}\le B \qquad\perogni n\in\N .
\end{equation}
L'intervallo $[A,B]$ contiene tutti i termini, e quindi termini per infiniti indici
(perché $\N$ non è limitato superiormente, \cref{teo:archimede}).

\emph{Costruzione degli intervalli.} Suddividiamo l'intervallo $[A,B]$ mediante il punto
di mezzo $C=(A+B)/2$ e consideriamo i due intervalli $[A,C]$, $[C,B]$. Per quanto
osservato sopra, uno almeno dei due contiene termini della successione per infiniti
indici; indichiamolo con $[A_{1},B_{1}]$ (se entrambi lo fanno, ne scegliamo
uno qualsiasi). Risulta
\begin{equation}\label{eq:bw-passo1}
  A\le A_{1}, \qquad B_{1}\le B, \qquad B_{1}-A_{1}=\frac{B-A}{2}.
\end{equation}
Suddividiamo l'intervallo $[A_{1},B_{1}]$ mediante il punto di mezzo
$C_{1}=(A_{1}+B_{1})/2$. Per lo stesso motivo almeno uno tra i due intervalli
$[A_{1},C_{1}]$, $[C_{1},B_{1}]$ contiene termini della successione per infiniti indici;
indichiamolo con $[A_{2},B_{2}]$:
\begin{equation}\label{eq:bw-passo2}
  A_{1}\le A_{2}, \qquad B_{2}\le B_{1}, \qquad B_{2}-A_{2}=\frac{B_{1}-A_{1}}{2}=\frac{B-A}{2^{2}}.
\end{equation}
Iterando il procedimento (cioè definendo $[A_{k+1},B_{k+1}]$ a partire da $[A_{k},B_{k}]$,
per induzione) si generano due successioni $A_{k}$, $B_{k}$ $(k\in\N)$ tali che
\begin{gather}
  A\le A_{k}\le A_{k+1}<B_{k+1}\le B_{k}\le B \qquad\perogni k\in\N, \label{eq:bw-incapsulati}\\
  B_{k}-A_{k}=\frac{B-A}{2^{k}} \qquad\perogni k\in\N, \label{eq:bw-lunghezza}
\end{gather}
e inoltre ogni intervallo $[A_{k},B_{k}]$ contiene termini della successione $a_{n}$ per
infiniti indici.

\emph{Scelta degli indici.} L'intervallo $[A_{1},B_{1}]$ contiene termini della
successione $a_{n}$; quindi esiste il primo intero $n_{1}$ tale che
$a_{n_{1}}\in[A_{1},B_{1}]$ (\cref{teo:buon-ordinamento}). Per lo stesso motivo esiste un
primo intero $n_{2}$, fra tutti i numeri naturali più grandi di $n_{1}$, per cui
$a_{n_{2}}\in[A_{2},B_{2}]$: ne esistono, perché $[A_{2},B_{2}]$ contiene termini per
infiniti indici. Iterando, per $k=3,4,5,\dots$ determiniamo una successione
strettamente crescente di interi
\begin{equation}\label{eq:bw-indici}
  n_{1}<n_{2}<n_{3}<\dots<n_{k}<n_{k+1}<\dots
\end{equation}
per cui $a_{n_{k}}\in[A_{k},B_{k}]$ per ogni $k\in\N$. Dato che
$B_{k}-A_{k}=(B-A)/2^{k}$, abbiamo quindi
\begin{equation}\label{eq:bw-carabinieri}
  A_{k}\le a_{n_{k}}\le B_{k}=A_{k}+\frac{B-A}{2^{k}} \qquad\perogni k\in\N .
\end{equation}

\emph{Passaggio al limite.} Per la \eqref{eq:bw-incapsulati} la successione $A_{k}$ è
monotòna (crescente) e limitata; per il teorema sulle successioni monotòne $A_{k}$ ammette
limite finito per $k\to+\infty$. Indichiamo con $\ell\in\R$ il valore di tale limite. Dato
che $(B-A)/2^{k}=(B-A)\bigl(\frac12\bigr)^{k}\to0$ per $k\to+\infty$
(\cref{teo:lim-geometrica}), sia il primo sia l'ultimo membro della
\eqref{eq:bw-carabinieri} convergono a $\ell$. Inoltre $A\le A_{k}\le B$ per ogni $k$, per
la \eqref{eq:bw-incapsulati}, e quindi $A\le\ell\le B$ per il \cref{cor:confronto-limiti}
(applicato confrontando $A_{k}$ con le successioni costanti $A$ e $B$): dunque
$\ell\in[A,B]$. Dal teorema dei carabinieri si ottiene allora la conclusione
\begin{equation}\label{eq:bw-tesi}
  \lim_{k\to+\infty}a_{n_{k}}=\ell. \qedhere
\end{equation}
\end{dimostrazione}

\begin{figure}[htb]
\centering
\begin{tikzpicture}[x=10cm, every node/.style={font=\footnotesize}]
  % termini di una successione limitata (valori di sen n riscalati in [0,1])
  \foreach \row/\A/\B in {0/0/1, 1/0.5/1, 2/0.5/0.75, 3/0.625/0.75, 4/0.625/0.6875} {
    \begin{scope}[yshift={-\row*0.85cm}]
      \draw[grigio!60] (0,0) -- (1,0);
      \fill[blu!18] (\A,-0.12) rectangle (\B,0.12);
      \draw[blu, thick] (\A,-0.12) -- (\A,0.12); \draw[blu, thick] (\B,-0.12) -- (\B,0.12);
      \ifnum\row=0 \node[left] at (0,0) {$[A,B]$}; \else \node[left] at (0,0) {$[A_{\row},B_{\row}]$}; \fi
    \end{scope}}
  \foreach \n in {1,...,60} {
    \pgfmathsetmacro{\v}{0.5+0.5*sin(deg(\n))}
    \fill[rame] (\v,0) circle (0.9pt);}
  \draw[verde!60!black, thick, dashed] (0.6553,0.3) -- (0.6553,-3.8) node[below] {$\ell$};
\end{tikzpicture}
\caption{Il procedimento di bisezione nella dimostrazione del teorema di
Bolzano--Weierstrass. I punti color rame sono i termini (riscalati) di una successione
limitata; a ogni passo si conserva una metà dell'intervallo che contiene termini per
infiniti indici. Gli intervalli si stringono attorno al limite $\ell$ dell'estratta.}
\label{fig:bolzano-weierstrass}
\end{figure}

\begin{osservazione}
Il teorema afferma l'esistenza di un'estratta convergente, ma non ne fornisce, in
generale, una costruzione effettiva: la scelta della metà «buona» a ogni passo richiede
di sapere quale delle due contiene infiniti termini, cosa che nessun calcolatore può
verificare. Inoltre l'estratta non è unica, e successioni diverse possono avere estratte
con limiti diversi: per $\sen n$, ad esempio, si può dimostrare che \emph{ogni} numero di
$[-1,1]$ è limite di qualche estratta. Infine, l'ipotesi di limitatezza è essenziale: la
successione $a_{n}=n$ non ha alcuna estratta convergente, perché ogni estratta diverge a
$+\infty$ (\cref{prop:estratta-limite}).
\end{osservazione}

\subsection{Successioni di Cauchy}

Il teorema sulle successioni monotòne permette di dimostrare che una successione converge
senza conoscerne il limite; ma richiede la monotonia. Il criterio di Cauchy offre un
analogo strumento per successioni qualsiasi.

\begin{definizione}[Successione di Cauchy]\index{successione!di Cauchy}
Sia $a_{n}$ una successione di numeri reali. Si dice che $a_{n}$ è una \emph{successione di
Cauchy} se, per ogni $\epsilon>0$, esiste un indice $\nu$ tale che per $h,k>\nu$ risulti
\begin{equation}\label{eq:cauchy}
  \abs{a_{k}-a_{h}}<\epsilon .
\end{equation}
\end{definizione}

La condizione di Cauchy non menziona alcun limite: chiede soltanto che i termini, da un
certo punto in poi, siano \emph{vicini tra loro}. Dimostriamo in primo luogo la seguente

\begin{proposizione}\label{prop:convergente-cauchy}
Ogni successione convergente è di Cauchy.
\end{proposizione}

\begin{schema}
\Ipotesi{$a_{n}\to a\in\R$.}
\Tesi{$\perogni\epsilon>0\ \esiste\nu:\ \abs{a_{k}-a_{h}}<\epsilon$ per ogni $h,k>\nu$.}
\end{schema}

\begin{strategia}
Se due termini sono entrambi vicini al limite $a$, sono vicini tra loro: la disuguaglianza
triangolare, passando per $a$, somma le due distanze. Per ottenere $\epsilon$ alla fine
si chiede a ciascuna distanza di essere minore di $\epsilon/2$.
\end{strategia}

\begin{dimostrazione}
Se $a_{n}$ converge verso $a$, allora, per ogni $\epsilon>0$, esiste $\nu$ tale che
\begin{equation}\label{eq:conv-cauchy-1}
  \abs{a_{n}-a}<\frac\epsilon2 \qquad\perogni n>\nu .
\end{equation}
Dalla disuguaglianza triangolare segue allora, per $h,k>\nu$:
\begin{equation}\label{eq:conv-cauchy-2}
  \abs{a_{k}-a_{h}}\le\abs{a_{k}-a}+\abs{a-a_{h}}<\frac\epsilon2+\frac\epsilon2=\epsilon. \qedhere
\end{equation}
\end{dimostrazione}

Per dimostrare che, viceversa, ogni successione di Cauchy è convergente, premettiamo alcuni
lemmi.

\begin{lemma}\label{lem:cauchy-limitata}
Una successione di Cauchy è limitata.
\end{lemma}

\begin{schema}
\Ipotesi{$a_{n}$ di Cauchy.}
\Tesi{$\esiste A,B\in\R$: \ $A\le a_{n}\le B$ per ogni $n\in\N$.}
\end{schema}

\begin{strategia}
Con $\epsilon=1$, da un certo indice in poi tutti i termini distano meno di $1$ da un
termine fissato $a_{h_{0}}$: sono quindi definitivamente in un intervallo limitato.
Restano fuori un numero finito di termini, che si aggiungono al calcolo del minimo e del
massimo.
\end{strategia}

\begin{dimostrazione}
Sia $\epsilon=1$; per ipotesi esiste $\nu\in\N$ tale che
\begin{equation}\label{eq:cauchy-lim-1}
  \abs{a_{k}-a_{h}}<1 \qquad\perogni h,k>\nu .
\end{equation}
Fissiamo un indice $h_{0}>\nu$. Allora, dalla \eqref{eq:cauchy-lim-1}, per le proprietà
del valore assoluto, segue
\begin{equation}\label{eq:cauchy-lim-2}
  a_{h_{0}}-1<a_{k}<a_{h_{0}}+1 \qquad\perogni k>\nu .
\end{equation}
Posto
\[
  A=\min\{a_{1},\dots,a_{\nu},a_{h_{0}}-1\}, \qquad B=\max\{a_{1},\dots,a_{\nu},a_{h_{0}}+1\},
\]
evidentemente risulta $A\le a_{k}\le B$ per ogni $k\in\N$, e perciò la successione è
limitata. (È lo stesso schema del \cref{teo:convergente-limitata}: controllo definitivo
più un numero finito di eccezioni.)
\end{dimostrazione}

\begin{lemma}\label{lem:cauchy-estratta}
Se una successione di Cauchy $a_{n}$ contiene un'estratta $a_{n_{k}}$ convergente verso
$\ell$, allora anche $a_{n}$ converge verso $\ell$.
\end{lemma}

\begin{schema}
\Ipotesi{$a_{n}$ di Cauchy; \ $a_{n_{k}}\to\ell\in\R$.}
\Tesi{$a_{n}\to\ell$.}
\end{schema}

\begin{strategia}
Per stimare $\abs{a_{n}-\ell}$ si passa per un termine dell'estratta abbastanza avanzato,
$a_{n_{k_{0}}}$: questo è vicino a $\ell$ perché l'estratta converge, ed è vicino ad
$a_{n}$ perché la successione è di Cauchy. La disuguaglianza triangolare, con il solito
$\epsilon/2+\epsilon/2$, conclude.
\end{strategia}

\begin{dimostrazione}
Fissato $\epsilon>0$, sia $\nu\in\N$ tale che
\begin{equation}\label{eq:cauchy-estr-1}
  \abs{a_{k}-a_{h}}<\frac\epsilon2 \qquad\perogni h,k>\nu .
\end{equation}
Sia inoltre $k_{0}>\nu$ tale che
\begin{equation}\label{eq:cauchy-estr-2}
  \abs{a_{n_{k}}-\ell}<\frac\epsilon2 \qquad\perogni k\ge k_{0}.
\end{equation}
Poiché si ha $n_{k_{0}}\ge k_{0}>\nu$ (\cref{lem:nk-maggiore-k}), per ogni $n>\nu$ risulta
\begin{equation}\label{eq:cauchy-estr-3}
  \abs{a_{n}-\ell}\le\abs{a_{n}-a_{n_{k_{0}}}}+\abs{a_{n_{k_{0}}}-\ell}<\frac\epsilon2+\frac\epsilon2=\epsilon. \qedhere
\end{equation}
\end{dimostrazione}

Combinando i risultati precedenti si dimostra il seguente

\begin{teorema}[Criterio di convergenza di Cauchy]\label{teo:criterio-cauchy}\index{criterio di Cauchy!per le successioni}
Una successione $a_{n}$ di numeri reali è convergente se e solo se è di Cauchy.
\end{teorema}

\begin{schema}
\Ipotesi{$a_{n}$ successione di numeri reali.}
\Tesi{$a_{n}$ converge in $\R$ \ $\sse$ \ $a_{n}$ è di Cauchy, cioè per ogni $\epsilon>0$
esiste $\nu$ tale che $\abs{a_{k}-a_{h}}<\epsilon$ per ogni $h,k>\nu$.}
\end{schema}

\begin{strategia}
Una delle due implicazioni è la \cref{prop:convergente-cauchy}. Per l'altra, il problema è
che non conosciamo il candidato limite; lo facciamo produrre al teorema di
Bolzano--Weierstrass. Una successione di Cauchy è limitata (\cref{lem:cauchy-limitata}),
quindi ha un'estratta convergente a qualche $\ell$; e poiché i termini di una successione
di Cauchy sono definitivamente vicini tra loro, se alcuni di essi si avvicinano a $\ell$,
vi si avvicinano tutti (\cref{lem:cauchy-estratta}).
\end{strategia}

\begin{dimostrazione}
Con la \cref{prop:convergente-cauchy} abbiamo già provato che le successioni convergenti
sono di Cauchy. Viceversa, se $a_{n}$ è di Cauchy, per il \cref{lem:cauchy-limitata} è
anche limitata. In base al teorema di Bolzano--Weierstrass (\cref{teo:bolzano-weierstrass})
$a_{n}$ ammette una successione estratta $a_{n_{k}}$ convergente. Per il
\cref{lem:cauchy-estratta}, $a_{n}$ è convergente.
\end{dimostrazione}

\begin{osservazione}[Cauchy e la completezza]
Nel campo dei razionali il criterio di Cauchy è falso: la successione dei troncamenti
decimali di $\sqrt2$,
\[
  1;\quad 1{,}4;\quad 1{,}41;\quad 1{,}414;\quad 1{,}4142;\quad\dots,
\]
è formata da numeri razionali ed è di Cauchy (due termini di indice maggiore di $k$
distano meno di $10^{-k}$), ma non converge ad alcun numero razionale, perché il suo
limite in $\R$ è $\sqrt2\notin\Q$. La dimostrazione del criterio usa infatti, attraverso
il teorema di Bolzano--Weierstrass e quello sulle successioni monotòne, l'assioma di
completezza. Si può anzi dimostrare che, in un campo ordinato archimedeo, il criterio di
Cauchy è \emph{equivalente} all'assioma di completezza \cite{giusti}. Lo strumento è
ancora la bisezione. Dati $A$ e $B$ come nell'assioma, si sceglie $a_{0}\in A$ e
$b_{0}\in B$ e si dimezza ripetutamente l'intervallo $[a_{0},b_{0}]$, conservando a ogni
passo una metà il cui estremo destro è un maggiorante di $A$ e il cui estremo sinistro non
supera qualche elemento di $A$. Gli estremi formano due successioni di Cauchy, e il loro
limite comune separa $A$ da $B$. La proprietà di Archimede serve a garantire che le
lunghezze $(b_{0}-a_{0})/2^{k}$ tendano a zero.
\end{osservazione}

\begin{python}[Una successione di Cauchy di razionali che non converge in $\Q$]
# Troncamenti decimali di sqrt(2): sono razionali, distano sempre
# meno, ma il loro limite sqrt(2) non è razionale.
from fractions import Fraction as F
from decimal import Decimal, getcontext
getcontext().prec = 40
cifre = str(Decimal(2).sqrt())            # "1.414213562373095..."

tronc = [F(cifre[:2 + k]) for k in range(1, 12)]   # 1.4, 1.41, ...
print("k               q_k   max |q_j - q_k|, j > k   q_k^2 = 2 ?")
for k in range(1, 8):
    q = tronc[k - 1]
    dist = max(abs(tronc[j] - q) for j in range(k, len(tronc)))
    print(f"{k}  {str(q):>16s}   {float(dist):22.1e}   {q * q == 2}")
\end{python}
\begin{risultato}
k               q_k   max |q_j - q_k|, j > k   q_k^2 = 2 ?
1               7/5                  1.4e-02   False
2           141/100                  4.2e-03   False
3           707/500                  2.1e-04   False
4         7071/5000                  1.4e-05   False
5     141421/100000                  3.6e-06   False
6   1414213/1000000                  5.6e-07   False
7   2828427/2000000                  6.2e-08   False
\end{risultato}

\begin{idea}[I numeri reali come completamento dei razionali]
Il programma mostra una successione di Cauchy di razionali che «vorrebbe» convergere ma
non trova in $\Q$ il suo limite. Charles Méray (1869) e, indipendentemente, Georg Cantor
(1872) rovesciarono il problema: se il limite manca, lo si \emph{definisce} a partire
dalla successione stessa. Due successioni di Cauchy di razionali $a_{n}$ e $b_{n}$ si
dicono equivalenti se $a_{n}-b_{n}\to0$: ad esempio i troncamenti decimali di $\sqrt2$ e
i termini $2,\frac32,\frac{17}{12},\frac{577}{408},\dots$ del metodo di Erone
(\cref{es:erone}). Un numero reale è, per definizione, una classe di successioni
equivalenti, e il razionale $q$ si identifica con la classe della successione costante
$q$. Somma e prodotto si definiscono termine a termine; un numero è $\ge0$ se ha un
rappresentante definitivamente $\ge0$. Con un lavoro lungo ma elementare (la parte più
delicata riguarda l'ordinamento) si verifica che l'insieme così costruito soddisfa tutti
gli assiomi del \cref{sec:assiomi}, compreso quello di completezza \cite{giusti}. Anche
$e$ trova posto: è la classe della successione di razionali $(1+\frac1n)^{n}$, alla
quale è equivalente quella delle somme $\sum_{k=0}^{n}\frac1{k!}$
(\cref{oss:e-due-vie}). È la costruzione seguita in \cite{tao1}, alternativa alle
sezioni di Dedekind ricordate nella nota storica del \S\ref{sec:assiomi}.
\end{idea}

\begin{notastorica}[Bolzano, Cauchy, Weierstrass]
L'idea della dimostrazione del teorema che chiamiamo «di Bolzano--Weierstrass», la
bisezione, compare già nel saggio di Bernard Bolzano del 1817 ricordato nel
\cref{sec:sup}, dove serve a dimostrare l'esistenza dell'estremo superiore e, da questa,
il teorema dei valori intermedi. Nello stesso lavoro Bolzano enunciò
anche, quattro anni prima di Cauchy, il criterio che porta il nome di quest'ultimo, e
cercò di dimostrarlo. Il suo argomento si riduceva però a osservare che l'ipotesi di un
limite «non contiene nulla di impossibile», perché, se il limite c'è, il suo valore si
può determinare con tutta la precisione che si vuole: così si descrive il limite, ma non
se ne prova l'esistenza. In mancanza di una teoria dei numeri reali la dimostrazione non
poteva che essere illusoria \cite{giusti}. Il saggio, d'altra parte, rimase quasi
sconosciuto, e fu Karl Weierstrass a riscoprire il risultato e a
farne, nelle lezioni berlinesi degli anni Sessanta dell'Ottocento, uno dei pilastri
dell'Analisi. Augustin-Louis Cauchy enunciò il suo criterio nel \emph{Cours d'analyse}
(1821), ma ne considerò ovvia la parte «difficile» (che una successione di Cauchy
converge): ovvia lo è, se si ha in mente la retta geometrica; dimostrarla richiede una
teoria dei numeri reali, che né a Cauchy né a Bolzano era disponibile. Fu proprio per
colmare questa lacuna che, cinquant'anni dopo, Méray, Dedekind e Cantor costruirono i
numeri reali (si veda il riquadro precedente).
\end{notastorica}

\subsection{Approfondimento: massimo e minimo limite}
\label{subsec:massimo-minimo-limite}
\index{massimo limite}\index{minimo limite}

Una successione limitata può non avere limite, come $(-1)^{n}$. Il teorema di
Bolzano--Weierstrass garantisce che abbia estratte convergenti, ma non dice quali, e i
loro limiti possono essere diversi. Vi sono però due numeri, associati a ogni
successione limitata, che esistono \emph{sempre} e che si riducono al limite quando
questo c'è: il \emph{massimo} e il \emph{minimo limite} \cite{giusti}. L'idea è guardare
la successione «da un certo punto in poi», cioè le sue \emph{code}, e misurarne il
«tetto» e il «pavimento».

Sia $a_{n}$ una successione limitata, diciamo $A\le a_{n}\le B$ per ogni $n$. Per ogni
$n\in\N$ la coda $\{a_{k} : k\ge n\}$ è un insieme non vuoto e limitato: per il
\cref{teo:esistenza-sup} ha estremo superiore ed estremo inferiore reali, che indichiamo
con
\begin{equation}\label{eq:code-limsup}
  s_{n}=\sup_{k\ge n}a_{k}, \qquad i_{n}=\inf_{k\ge n}a_{k} .
\end{equation}
Poiché $a_{n}$ appartiene alla propria coda, si ha $A\le i_{n}\le a_{n}\le s_{n}\le B$.
Passando da $n$ a $n+1$, inoltre, la coda perde il termine $a_{n}$. Il numero $s_{n}$,
che maggiora la coda di indice $n$, maggiora a maggior ragione quella di indice $n+1$;
e poiché $s_{n+1}$ è il \emph{minimo} dei maggioranti di quest'ultima, $s_{n+1}\le s_{n}$.
Allo stesso modo $i_{n+1}\ge i_{n}$. La successione $s_{n}$ è dunque decrescente, la
$i_{n}$ crescente, ed entrambe sono limitate: per il \cref{teo:monotone} convergono.

\begin{definizione}[Massimo e minimo limite]\label{def:limsup}
Sia $a_{n}$ una successione limitata, e siano $s_{n}$ e $i_{n}$ definite dalla
\eqref{eq:code-limsup}. I numeri reali
\begin{equation}\label{eq:limsup-def}
\begin{aligned}
  \limsup_{n\to+\infty}a_{n}&=\lim_{n\to+\infty}s_{n}=\inf_{n\in\N}s_{n},\\
  \liminf_{n\to+\infty}a_{n}&=\lim_{n\to+\infty}i_{n}=\sup_{n\in\N}i_{n}
\end{aligned}
\end{equation}
si chiamano, rispettivamente, \emph{massimo limite} e \emph{minimo limite} (o
\emph{limite superiore} e \emph{limite inferiore}) della successione $a_{n}$. Si usano
anche i simboli $\overline{\lim}\,a_{n}$ e $\underline{\lim}\,a_{n}$.
\end{definizione}

Il massimo limite è dunque il tetto della successione «visto da lontano»: $s_{n}$ è il
più basso dei livelli che i termini da $a_{n}$ in poi non superano, e
$\limsup a_{n}$ è il valore verso cui questi livelli si abbassano al crescere di $n$.

\begin{esempio}[Una successione alternante]\label{es:limsup-alternante}
Sia $a_{n}=(-1)^{n}\bigl(1+\frac1n\bigr)$. I termini di indice pari, $1+\frac1n$,
decrescono verso $1$; quelli di indice dispari, $-1-\frac1n$, crescono verso $-1$
(\cref{fig:limsup-code}). Il più grande termine della coda di indice $n$ è il primo
termine di indice pari che vi compare: $s_{n}=1+\frac1n$ se $n$ è pari,
$s_{n}=1+\frac1{n+1}$ se $n$ è dispari. Analogamente $i_{n}=-1-\frac1n$ se $n$ è dispari
e $i_{n}=-1-\frac1{n+1}$ se $n$ è pari. Quindi
\[
  \limsup_{n\to+\infty}a_{n}=1, \qquad \liminf_{n\to+\infty}a_{n}=-1,
\]
mentre $\sup_{n}a_{n}=a_{2}=\frac32$ e $\inf_{n}a_{n}=a_{1}=-2$. Gli estremi
superiore e inferiore dipendono dai primi termini; il massimo e il minimo limite no,
perché modificando un numero finito di termini le code da un certo indice in poi non
cambiano.
\end{esempio}

\begin{figure}[htb]
\centering
\begin{tikzpicture}
\begin{axis}[width=0.92\textwidth, height=5.6cm, xmin=0, xmax=25, ymin=-2.25, ymax=1.85,
  xtick={5,10,15,20}, ytick={-2,-1,1,1.5}, yticklabels={$-2$,$\lambda=-1$,$L=1$,$\frac32$},
  xlabel={$n$}, legend columns=3,
  legend style={font=\scriptsize, draw=none, fill=none, at={(0.5,0.69)}, anchor=center}]
  \addplot[tratteggio, domain=0:25, forget plot] {1};
  \addplot[tratteggio, domain=0:25, forget plot] {-1};
  \addplot[only marks, mark=*, mark size=1.4pt, blu, samples at={1,...,24}] {(-1)^x*(1+1/x)};
  \addplot[rame, thick, const plot, samples at={1,...,25}] {1+1/(2*ceil(x/2))};
  \addplot[verde, thick, const plot, samples at={1,...,25}] {-1-1/(2*floor(x/2)+1)};
  \legend{$a_{n}$, $s_{n}=\sup_{k\ge n}a_{k}$, $i_{n}=\inf_{k\ge n}a_{k}$}
\end{axis}
\end{tikzpicture}
\caption{La successione $a_{n}=(-1)^{n}\bigl(1+\frac1n\bigr)$
dell'\cref{es:limsup-alternante}, con i «tetti» $s_{n}$ (in rame) e i «pavimenti»
$i_{n}$ (in verde) delle sue code, disegnati a gradini. I tetti scendono verso
$L=\limsup a_{n}=1$, i pavimenti salgono verso $\lambda=\liminf a_{n}=-1$; l'estremo
superiore $\frac32$ e l'estremo inferiore $-2$ sono invece fissati dai primi termini.}
\label{fig:limsup-code}
\end{figure}

\begin{idea}[Il picco a regime di una risposta oscillante]
In una prova di vibrazione, o nella simulazione della risposta di un'ala a una raffica
periodica, si registra a istanti successivi una grandezza, ad esempio lo spostamento
della punta alare. La successione dei campioni contiene un \emph{transitorio}, spesso
con un picco iniziale elevato, seguito da un'oscillazione a regime. L'estremo superiore
dei campioni risente del picco iniziale; il massimo limite no: misura il picco \emph{a
regime}, quello che si ripete a ogni ciclo e che conta, ad esempio, per la fatica del
materiale. Allo stesso modo il minimo limite misura la valle a regime, e la semidifferenza
$\frac12(L-\lambda)$ l'ampiezza dell'oscillazione stabilizzata. Nell'\cref{es:limsup-alternante}
il valore $\sup a_{n}=\frac32$ è dovuto al secondo termine, mentre a regime la
successione oscilla tra valori sempre più vicini a $-1$ e a $1$.
\end{idea}

\begin{proposizione}[Proprietà del massimo e del minimo limite]\label{prop:limsup-proprieta}
Sia $a_{n}$ una successione limitata, e siano $L=\limsup a_{n}$ e $\lambda=\liminf a_{n}$.
Allora:
\begin{enumerate}
  \item $\inf_{n}a_{n}\le\lambda\le L\le\sup_{n}a_{n}$;
  \item per ogni $\epsilon>0$ si ha $a_{n}<L+\epsilon$ definitivamente, e
  $a_{n}>L-\epsilon$ per infiniti indici;
  \item per ogni $\epsilon>0$ si ha $a_{n}>\lambda-\epsilon$ definitivamente, e
  $a_{n}<\lambda+\epsilon$ per infiniti indici.
\end{enumerate}
\end{proposizione}

\begin{schema}
\Ipotesi{$a_{n}$ limitata; \ $L=\limsup a_{n}$, \ $\lambda=\liminf a_{n}$; \ $\epsilon>0$.}
\Tesi{(i) $\inf a_{n}\le\lambda\le L\le\sup a_{n}$; \ (ii) $a_{n}<L+\epsilon$
definitivamente, e $a_{n}>L-\epsilon$ per infiniti indici; \ (iii) le affermazioni
simmetriche per $\lambda$.}
\end{schema}

\begin{strategia}
La (i) si ottiene passando al limite nella catena $i_{1}\le i_{n}\le s_{n}\le s_{1}$. La
(ii) traduce le caratterizzazioni dell'estremo inferiore e dell'estremo superiore.
Poiché $L$ è l'estremo inferiore dei tetti $s_{n}$, qualche tetto $s_{\nu}$ è più basso
di $L+\epsilon$, e da $\nu$ in poi i termini non lo superano. D'altra parte ogni tetto
$s_{m}$ è almeno $L$: quindi in ogni coda qualche termine supera $L-\epsilon$, e ciò
accade oltre qualunque indice $m$. La (iii) è simmetrica.
\end{strategia}

\begin{dimostrazione}
(i) Per ogni $n$ si ha $i_{1}\le i_{n}\le s_{n}\le s_{1}$, dove $i_{1}=\inf_{k}a_{k}$ e
$s_{1}=\sup_{k}a_{k}$ sono gli estremi di tutta la successione. Passando al limite per
$n\to+\infty$ (\cref{cor:confronto-limiti}, con il confronto anche con le successioni
costanti $i_{1}$ e $s_{1}$) si ottiene la (i).

(ii) Sia $\epsilon>0$. Poiché $L=\inf_{n}s_{n}$, per la \eqref{eq:caratt-inf} esiste
$\nu\in\N$ con $s_{\nu}<L+\epsilon$. Per ogni $n>\nu$ il termine $a_{n}$ appartiene alla
coda di indice $\nu$, e quindi $a_{n}\le s_{\nu}<L+\epsilon$: la prima affermazione è
provata. Per la seconda fissiamo $m\in\N$. Poiché $L$ è un minorante dei tetti,
$s_{m+1}\ge L>L-\epsilon$; per la caratterizzazione dell'estremo superiore
(\cref{prop:caratt-sup}) applicata alla coda di indice $m+1$, esiste $k>m$ con
$a_{k}>L-\epsilon$. Dato che $m$ è arbitrario, l'insieme degli indici $k$ con
$a_{k}>L-\epsilon$ non è limitato superiormente: la disuguaglianza vale per infiniti
indici.

(iii) Si ragiona allo stesso modo, scambiando i ruoli dell'estremo superiore e
dell'estremo inferiore.
\end{dimostrazione}

\begin{osservazione}[La definizione di \cite{giusti}]\label{oss:limsup-giusti}
Diciamo \emph{maggiorante definitivo} della successione un numero $M$ tale che
$a_{n}\le M$ definitivamente. La (ii) dice che ogni $M>L$ è un maggiorante definitivo
(con $\epsilon=M-L$), mentre nessun $M<L$ lo è (con $\epsilon=L-M$ si ha $a_{n}>M$ per
infiniti indici). Quindi $L$ è un minorante dell'insieme dei maggioranti definitivi, e
$L+\epsilon$ vi appartiene per ogni $\epsilon>0$: per la \eqref{eq:caratt-inf}, $L$ ne è
l'estremo inferiore. È la definizione di massimo limite adottata in \cite{giusti}, che è
dunque equivalente alla nostra; allo stesso modo il minimo limite è l'estremo superiore
dei minoranti definitivi. Lo stesso ragionamento mostra che le due condizioni della (ii)
\emph{caratterizzano} $L$: un numero che le soddisfa entrambe è l'estremo inferiore dei
maggioranti definitivi, e quindi coincide con $L$.
\end{osservazione}

\begin{teorema}[Massimo e minimo limite ed estratte]\label{teo:limsup-estratte}
Sia $a_{n}$ una successione limitata, e siano $L=\limsup a_{n}$ e $\lambda=\liminf a_{n}$.
\begin{enumerate}
  \item Esistono un'estratta di $a_{n}$ che converge a $L$ e un'estratta che converge
  a $\lambda$.
  \item Se un'estratta $a_{n_{k}}$ converge a un numero $\ell$, allora
  $\lambda\le\ell\le L$.
\end{enumerate}
In altre parole, $L$ e $\lambda$ sono il più grande e il più piccolo dei limiti delle
estratte convergenti di $a_{n}$.
\end{teorema}

\begin{schema}
\Ipotesi{$a_{n}$ limitata; \ $L=\limsup a_{n}$, \ $\lambda=\liminf a_{n}$.}
\Tesi{(i) esistono estratte $a_{n_{k}}\to L$ e $a_{m_{k}}\to\lambda$; \ (ii) se
$a_{n_{k}}\to\ell$, allora $\lambda\le\ell\le L$.}
\end{schema}

\begin{strategia}
Per la (i) costruiamo l'estratta un termine alla volta, come nella dimostrazione del
teorema di Bolzano--Weierstrass, ma senza bisezione. Per la
\cref{prop:limsup-proprieta}, per ogni $k$ gli indici $n$ con
$L-\frac1k<a_{n}<L+\frac1k$ sono infiniti, e se ne può sempre scegliere uno maggiore del
precedente. Per la (ii) basta osservare che $a_{n_{k}}$, avendo indice $n_{k}\ge k$, sta
nella coda di indice $k$: è quindi compreso tra $i_{k}$ e $s_{k}$, e si passa al limite.
\end{strategia}

\begin{dimostrazione}
(i) Sia $k\in\N$. Per la \cref{prop:limsup-proprieta}(ii) con $\epsilon=\frac1k$ esiste
$\nu_{k}$ tale che $a_{n}<L+\frac1k$ per ogni $n>\nu_{k}$, mentre $a_{n}>L-\frac1k$ per
infiniti indici, cioè per indici maggiori di qualunque numero fissato. Ne segue che, per
ogni $m$, esiste $n>\max\{m,\nu_{k}\}$ tale che
\begin{equation}\label{eq:limsup-scelta}
  L-\frac1k<a_{n}<L+\frac1k .
\end{equation}
Scegliamo allora un indice $n_{1}$ che soddisfi la \eqref{eq:limsup-scelta} con $k=1$,
poi un indice $n_{2}>n_{1}$ che la soddisfi con $k=2$, e così via. Per induzione si
ottiene una successione strettamente crescente di indici $n_{k}$ con
$L-\frac1k<a_{n_{k}}<L+\frac1k$ per ogni $k$. Poiché $L\pm\frac1k\to L$, per il teorema
dei carabinieri $a_{n_{k}}\to L$. In modo analogo, usando la (iii), si costruisce
un'estratta convergente a $\lambda$.

(ii) Per il \cref{lem:nk-maggiore-k} si ha $n_{k}\ge k$, quindi $a_{n_{k}}$ appartiene
alla coda di indice $k$, e
\[
  i_{k}\le a_{n_{k}}\le s_{k} \qquad\perogni k\in\N .
\]
Per $k\to+\infty$ si ha $i_{k}\to\lambda$, $s_{k}\to L$ e $a_{n_{k}}\to\ell$; per il
\cref{cor:confronto-limiti} risulta $\lambda\le\ell\le L$.
\end{dimostrazione}

\begin{osservazione}[Un'altra dimostrazione del teorema di Bolzano--Weierstrass]\label{oss:bw-limsup}
Il punto (i) del \cref{teo:limsup-estratte} contiene una nuova dimostrazione del
\cref{teo:bolzano-weierstrass}: una successione limitata ha un'estratta convergente, ad
esempio quella che converge al massimo limite. È la via seguita da \cite{giusti}, e non è
circolare: abbiamo usato l'esistenza degli estremi superiore e inferiore, il teorema
sulle successioni monotòne e la caratterizzazione dell'estremo superiore, non il teorema
di Bolzano--Weierstrass. Il confronto tra le due vie è istruttivo. La bisezione cattura
un punto attorno al quale si addensano infiniti termini, senza dire quale; i tetti
$s_{n}$ indicano invece un punto preciso, il più grande possibile. In entrambe la
completezza entra nello stesso modo, attraverso il limite di una successione monotòna:
gli estremi $A_{k}$ degli intervalli dimezzati nella prima, i tetti $s_{n}$ nella
seconda.
\end{osservazione}

\begin{teorema}[Limite, massimo e minimo limite]\label{teo:limsup-limite}
Sia $a_{n}$ una successione limitata e sia $\ell\in\R$. Allora
\begin{equation}\label{eq:limsup-limite}
  \lim_{n\to+\infty}a_{n}=\ell \quad\sse\quad
  \limsup_{n\to+\infty}a_{n}=\liminf_{n\to+\infty}a_{n}=\ell .
\end{equation}
\end{teorema}

\begin{schema}
\Ipotesi{$a_{n}$ limitata; \ $\ell\in\R$.}
\Tesi{$a_{n}\to\ell$ \ $\sse$ \ $\limsup a_{n}=\liminf a_{n}=\ell$.}
\end{schema}

\begin{strategia}
Se la successione converge, ogni sua estratta converge allo stesso limite; ma per il
teorema precedente ci sono estratte che tendono a $L$ e a $\lambda$, che sono quindi
uguali a $\ell$. Viceversa, ogni termine sta tra il pavimento e il tetto della propria
coda, $i_{n}\le a_{n}\le s_{n}$: se pavimento e tetto tendono allo stesso limite, i
carabinieri concludono.
\end{strategia}

\begin{dimostrazione}
($\implica$) Se $a_{n}\to\ell$, per la \cref{prop:estratta-limite} ogni estratta di
$a_{n}$ tende a $\ell$. Per il \cref{teo:limsup-estratte}(i) esiste un'estratta che
tende a $L=\limsup a_{n}$, e per l'unicità del limite (\cref{teo:unicita-limite}) è
$L=\ell$. Allo stesso modo $\lambda=\liminf a_{n}=\ell$.

($\Leftarrow$) Per ogni $n$ si ha $i_{n}\le a_{n}\le s_{n}$, e per ipotesi
$i_{n}\to\lambda=\ell$ e $s_{n}\to L=\ell$. Per il teorema dei carabinieri
(\cref{teo:carabinieri}) $a_{n}\to\ell$.
\end{dimostrazione}

\begin{osservazione}\label{oss:limsup-illimitate}
Poiché ogni successione convergente è limitata (\cref{teo:convergente-limitata}), il
teorema si legge anche così: \emph{una successione converge se e solo se è limitata e
il suo massimo limite coincide con il minimo limite}. La differenza $L-\lambda\ge0$
misura quanto la successione «oscilla» all'infinito. Per la successione
dell'\cref{es:limsup-alternante} vale $2$, e la successione non converge; per
$(-1)^{n}/n$ si ha $0<s_{n}\le\frac1n$ e $-\frac1n\le i_{n}<0$, quindi $L=\lambda=0$, che
è il limite. Con le convenzioni $\sup=+\infty$ e $\inf=-\infty$ per gli insiemi non
limitati, le definizioni si estendono a successioni qualsiasi. Ad esempio, per
$a_{n}=(-1)^{n}n$ ogni coda è illimitata sia superiormente sia inferiormente, e si pone
$\limsup a_{n}=+\infty$, $\liminf a_{n}=-\infty$. In questo senso il massimo e il minimo
limite esistono sempre, anche quando il limite non esiste.
\end{osservazione}
